\lecture{10}
\title{Computation History Method}

\begin{document}

\subsection{Computation History}

\textbf{Definition.}
The \textbf{configuration} $C$ of a TM $M$
refers to the state of its tape, its head position, and its finite control.

\textbf{Definition.}
A \textbf{computation history} for a TM $M$ on an input string $w$ is
a sequence of configurations $(C_1, C_2, \dots, C_{\ell})$ such that:
\begin{itemize}
    \item We have that $C_1$ is the start configuration of $M$ on $w$.
    \item Each $C_{i + 1}$ legally follows from $C_i$
    according to the rules of $M$.
\end{itemize}

\textbf{Definition.}
We say $(C_1, C_2, \dots, C_{\ell})$ is \textbf{accepting}
if $C_{\ell}$ is an accepting configuration, and \textbf{rejecting}
if $C_{\ell}$ is rejecting.

\textbf{Definition.}
We often encode configurations and computation histories
like so.
\begin{itemize}
    \item Let $\texttt{c}_i$ denote
    the encoding of a configuration $C_i$.
    \item Let the encoding of a computation history
    $H = (C_1, C_2, \dots, C_{\ell})$ be
    $\langle H \rangle =
    \texttt{c}_1 \texttt{\#}
    \texttt{c}_2 \texttt{\#}
    \dots
    \texttt{\#} \texttt{c}_{\ell}$.
\end{itemize}

\subsection{Linearly Bounded Automata:
$A_{\mathrm{LBA}}$}

\textbf{Definition.}
A \textbf{linearly bounded automaton (LBA)} is a Turing machine
in which the tape head is not permitted to move
off the portion of the tape containing the input---as in,
it only has as much memory as its input length.

\begin{theorem}{$A_{\mathrm{LBA}}$ Decidable}
    The language $A_{\mathrm{LBA}} := \{
        \langle B, w \rangle \mid
        B \text{ is an LBA that accepts } w
    \}$ is decidable.
\end{theorem}

\begin{proof}
    If $B$ is an LBA with $q$ states,
    a tape of length $n$,
    and $g$ letters in its tape alphabet,
    then $B$ only has $qng^n$ distinct configurations.
    So it suffices to simulate $B$ on $w$
    only for $qng^n$ steps.
    \begin{itemize}
        \item Return \textsc{Accept}
        if the simulation ever returns \textsc{Accept}.
        \item Return \textsc{Reject}
        if the simulation ever returns \textsc{Reject}
        or has not halted within $qng^n$ steps.
    \end{itemize}
    Indeed, if the simulation has not halted within $qng^n$ steps,
    it must loop and therefore never halt. ~ $\blacksquare$
\end{proof}

\subsection{The Computation History Method:
$E_{\mathrm{LBA}}$, $\mathrm{ALL}_{\mathrm{PDA}}$,
$E_{\mathrm{PDA}}$}

The value of computation histories becomes especially apparent
when we apply them to weaker models of computation than a Turing machine,
such as the LBA.

\begin{theorem}{$E_{\mathrm{LBA}}$ Undecidable}
    The language $E_{\mathrm{LBA}} := \{
        \langle B \rangle \mid
        B \text{ is an LBA such that } L(B) = \emptyset
    \}$ is undecidable.
\end{theorem}

\begin{proof}
    We reduce $A_{\mathrm{TM}}$
    (from \href{/18.404/lectures/8/#a-tm-undecidability}{Lecture 8})
    to $E_{\mathrm{LBA}}$.
    Given a TM $R$ that decides $E_{\mathrm{LBA}}$,
    we construct a TM $S$ that decides $A_{\mathrm{TM}}$.
    \begin{quote}
        \textbf{Claim.}
        We can construct an LBA $B_{M, w}$ satisfying the following.
        \begin{enumerate}
            \item The LBA $B_{M, w}$ takes a computation history
            $\langle H \rangle = \langle C_1, C_2, \dots, C_{\ell} \rangle$
            as input.
            \item The LBA $B_{M, w}$ accepts its input if and only if
            it is an accepting computation history for $M$ on $w$.
        \end{enumerate}

        \emph{Proof:}
        It suffices to verify each of the following.
        \begin{enumerate}
            \item The first configuration $C_1$
            agrees with taking $M$ on input $w$.
            \item Each configuration $C_{i + 1}$
            legally follows from the previous configuration $C_i$.
            \item The final configuration $C_{\ell}$
            is accepting.
        \end{enumerate}
        Each of these can easily be done by an LBA.
        ~ $\square$
    \end{quote}
    And so the TM $S$ that decides $A_{\mathrm{TM}}$
    can just run $R$ on the LBA $B_{M, w}$
    and output the opposite of what $R$ outputs.
    ~ $\blacksquare$
\end{proof}

Notice how an LBA is able to verify a computation history
even though it is unable to simulate a Turing machine.

Consequently, the utility of the computation history method
comes from its ability to make reductions from $A_{\mathrm{TM}}$
into languages concerning \emph{weaker} models of computation
than a Turing machine (e.g., an LBA).

The next example is similar---except instead of an LBA,
our ``\emph{weaker} model of computation'' will be a PDA.

\begin{theorem}{$\mathrm{ALL}_{\mathrm{PDA}}$ Undecidable}
    The language $\mathrm{ALL}_{\mathrm{PDA}} := \{
        \langle M \rangle \mid M \text{ is a PDA and } L(M) = \Sigma^*
    \}$ is undecidable.
\end{theorem}

\begin{proof}
    Reduce $A_{\mathrm{TM}}$ to $\mathrm{ALL}_{\mathrm{PDA}}$.
    Given a TM $R$ that decides $\mathrm{ALL}_{\mathrm{PDA}}$,
    we construct a TM $S$ that decides $A_{\mathrm{TM}}$.
    \begin{quote}
        \textbf{Claim.}
        We can construct a PDA $P_{M, w}$ satisfying the following.
        \begin{enumerate}
            \item The PDA $P_{M, w}$ takes a computation history
            $\langle H \rangle = \langle C_1, C_2, \dots, C_{\ell} \rangle$
            as input.
            \item The PDA $P_{M, w}$ accepts its input if and only if
            it is \textbf{not} an accepting computation history for $M$ on $w$.
        \end{enumerate}

        \emph{Proof:}
        A single PDA is able to perform
        any \emph{one} of the following checks.
        \begin{enumerate}
            \item The first configuration $C_1$
            does \emph{not} agree with taking $M$ on input $w$.
            \begin{itemize}
                \item To do this, $P_{M, w}$ can just have
                the correct first configuration hardcoded
                into one of its states.
            \end{itemize}
            \item For some fixed $i$, the configuration $C_{i + 1}$
            does \emph{not} legally follow from
            the previous configuration $C_i$.
            \begin{itemize}
                \item To do this,
                we actually need to slightly alter
                our encoding of a computation history
                to look like the following:
                \[
                \langle H \rangle =
                \texttt{c}_1\texttt{\#}
                \texttt{c}_2^R\texttt{\#}
                \texttt{c}_3\texttt{\#}
                \texttt{c}_4^R\texttt{\#}\dots
                ~~~ \text{ where } ~~~
                \langle C_i \rangle = \texttt{c}_i
                \text{ for all } i = 1, 2, \dots, \ell.
                \]
                This way, we can push $\texttt{c}_i$ onto the stack,
                then pop the stack to compare
                $\texttt{c}_i$ with $\texttt{c}_{i + 1}$
                character-by-character.
            \end{itemize}
            \item The final configuration $C_{\ell}$
            is \emph{not} accepting.
            \begin{itemize}
                \item To do this, $P_{M, w}$
                just needs to check whether $q_{\mathrm{acc}}$
                is in $C_{\ell}$.
            \end{itemize}
        \end{enumerate}
        The finish is to use the nondeterminism of PDAs
        to compute the logical OR of all of these checks---accept
        if \emph{any} of these checks are satisfied.
        ~ $\square$
    \end{quote}
    Then the TM $S$ that decides $A_{\mathrm{TM}}$
    can just run $R$ on the PDA $P_{M, w}$
    and output the opposite of what $R$ outputs.
    ~ $\blacksquare$
\end{proof}

\begin{remark}
    Recall from
    \href{/18.404/lectures/7/#algorithms-on-grammars-and-turing-machines-a-cfg-e-cfg-eq-cfg-a-tm}{Lecture 7}
    that the language $E_{\mathrm{PDA}} := \{
        \langle M \rangle \mid M \text{ is a PDA and }
        L(M) = \emptyset
    \}$ (equivalently, $E_{\mathrm{CFG}}$)
    actually \emph{is} decidable,
    unlike $\mathrm{ALL}_{\mathrm{PDA}}$.
    If one attempts to use the above argument
    to prove the undecidability of $E_{\mathrm{PDA}}$,
    this happens:
    \begin{quote}
        \textbf{Problem.}
        The PDA $P_{M, w}$ must accept its input
        if and only if it is an accepting computation history\dots

        \dots but PDAs aren't capable of verifying
        that a computation history is accepting!
    \end{quote}
    In particular, the nondeterminism of a PDA
    lets us compute the logical OR of multiple checks,
    but not the logical AND.
\end{remark}

\subsection{The Post Correspondence Problem}

We conclude by using computation histories
to prove a somewhat surprising result.

\begin{theorem}{PCP Undecidable}
    Consider the \textbf{Post Correspondence Problem (PCP)},
    defined as follows.
    \begin{quote}
        \textbf{Post Correspondence Problem (PCP).}
        Take a collection of \emph{dominos} as input $\langle P \rangle$.
        \[
        P = \left\{
            \left[
                \frac{t_1}{b_1}
            \right],
            \left[
                \frac{t_2}{b_2}
            \right],
            \dots,
            \left[
                \frac{t_k}{b_k}
            \right]
        \right\}
        ~~~~~~~~~~~ \textit{Example: }
        \left\{
            \left[
                \frac{\str{b}}{\str{ca}}
            \right],
            \left[
                \frac{\str{a}}{\str{ab}}
            \right],
            \left[
                \frac{\str{ca}}{\str{a}}
            \right],
            \left[
                \frac{\str{abc}}{\str{c}}
            \right]
        \right\}
        \]
        Accept if and only if the dominos may be arranged
        (possibly with repeats) such that
        the top and bottom rows match.
        \[
            \textit{Example: }
            \left[
                \frac{\str{a}}{\str{ab}}
            \right]
            \left[
                \frac{\str{b}}{\str{ca}}
            \right]
            \left[
                \frac{\str{ca}}{\str{a}}
            \right]
            \left[
                \frac{\str{a}}{\str{ab}}
            \right]
            \left[
                \frac{\str{abc}}{\str{c}}
            \right]
        \]
    \end{quote}
    It turns out, perhaps surprisingly, that PCP is \emph{undecidable}.
\end{theorem}



\begin{proof}
    Reduce $A_{\mathrm{TM}}$ to the PCP as before.
    First, a simplifying claim that we will assume without proof.

    \begin{quote}
        \textbf{Claim.}
        The \textbf{MPCP (Modified Post Correspondence Problem)}
        is a variant of the PCP in which the first domino is fixed.

        Then the PCP is decidable if and only if the MPCP is decidable.
    \end{quote}

    The idea is that, on input $\langle M, w \rangle$,
    we construct an MPCP instance such that
    the following holds.
    
    \begin{quote}
        \textbf{Goal.}
        Suppose there exists an arrangement of dominos
        in which the top and bottom rows both read $S$.
        Then $S$ must be the encoding of an accepting computation history
        for $M$ on $w$.
    \end{quote}

    Such a goal is pretty hard to achieve!
    In particular, the construction is extremely detail-heavy,
    so we won't fully justify it here.
    The interactive (warning: AI-generated) widget below
    should give a rough sense of how it works.

    \includewidget{lecture-10/mpcp}

    And so $A_{\mathrm{TM}}$ can be reduced to the MPCP
    by running an MPCP decider on
    the MPCP instance corresponding to $\langle M, w \rangle$.
    ~ $\blacksquare$
\end{proof}


\end{document}